\chapter{Introduction}
\label{ch:introduction}

Decision making under uncertainty is a central challenge in both theoretical and applied contexts, ranging from financial markets to artificial intelligence. Understanding how agents adapt their strategies when faced with noisy or incomplete information provides valuable insight into how learning and optimization unfold in dynamic environments. In particular, strategies that rely on probabilistic reasoning, such as Bayesian updating, the Kelly criterion, and multi-armed bandits, offer different ways of balancing risk and reward.

This thesis focuses on simulating a stylized market-like environment in which agents must repeatedly allocate resources across uncertain options. By examining how agents perform under varying conditions—including noisy feedback, limited horizons, and adversarial settings—this research evaluates the fundamental trade-offs between short-term performance and long-term optimality. 

\section{Research Objectives}
The primary objective is to develop a simulation framework that acts as a testbed for sequential decision-making strategies. We evaluate:
\begin{itemize}
    \item How the exploration versus exploitation dilemma manifests in continuous betting fraction spaces.
    \item The trade-off between maximizing long-term wealth versus maximizing the probability of reaching a short-term target while avoiding ruin.
    \item The robustness (or fragility) of standard Bayesian updating frameworks when faced with non-stationary, changing payoffs.
\end{itemize}

\section{Methodological Overview}
This work bridges Bayesian statistical modeling with probability and stochastic processes. Under the guidance of Professor Shiwei Lan, the project emphasizes Bayesian inference via algorithms like Thompson Sampling. Professor Nicolas Lanchier's expertise informs the probabilistic foundations of the Markov environments governing the hidden payoffs. By comparing experimental results with theoretical expectations, this thesis aims to construct a comprehensive understanding of adaptive probabilistic decision systems.
